\documentclass{ctexart}
\usepackage{amsmath, amssymb}
\usepackage{geometry}
\geometry{a4paper, margin=2cm}

\title{\textbf{第3章 功和能} --- 例题}
\author{}
\date{}

\begin{document}
\maketitle

\begin{enumerate}

\item 已知 $m = 2\,\text{kg}$，在 $F = 12t$ 作用下由静止做直线运动。求 $t = 0 \to 2\,\text{s}$ 内 $F$ 作的功及 $t = 2\,\text{s}$ 时的功率。

\item 质量为 $10\,\text{kg}$ 的质点，在外力作用下做平面曲线运动，该质点的速度为 $\vec{v} = 4t^{2}\,\vec{i} + 16\,\vec{j}$，开始时质点位于坐标原点。求在质点从 $y = 16\,\text{m}$ 到 $y = 32\,\text{m}$ 的过程中，外力做的功。

\item 已知用力 $\vec{F}$ 缓慢拉质量为 $m$ 的小球，$\vec{F}$ 保持方向不变。求 $\theta = \theta_0$ 时 $\vec{F}$ 作的功。

\item（书中例题 3.2，p.98，重点）一条长 $L$、质量 $M$ 的均匀柔绳，$A$ 端挂在天花板上，自然下垂，将 $B$ 端沿铅直方向提高到与 $A$ 端同高处。求该过程中重力所作的功。

\item（重点）蓄水池面积 $S$，水深 $h$，水面距地面 $H$，水的密度为 $\rho$。求抽出水需要作多少功？

\item 风力 $F$ 作用于向北运动的船，风力方向变化的规律是 $\theta = Bs$，其中 $s$ 为位移，$B$ 为常数，$\theta$ 为 $F$ 与 $S$ 间的夹角。如果运动中，风的方向自南变到东，求风力作的功。

\item（书中例题 3.5，p.103）物体的质量为 $m$，弹簧劲度系数为 $k$，$A$ 板及弹簧质量不计。求自弹簧原长 $O$ 处，突然无初速度地加上物体时，弹簧的最大压缩量。

\item（书中例题 3.11，p.111，重点）长为 $l$ 的均质链条，部分置于水平面上，另一部分自然下垂，已知链条与水平面间静摩擦系数为 $\mu_0$，滑动摩擦系数为 $\mu$。
\begin{enumerate}
  \item 满足什么条件时，链条将开始滑动？
  \item 若下垂部分长度为 $b$ 时，链条自静止开始滑动，当链条末端刚刚滑离桌面时，其速度等于多少？
\end{enumerate}

\item（书中例题 3.12，p.112）水平面内有一半径为 $R$ 的圆，在圆内离圆心 $O$ 距离为 $S$ 处有一质量很大、可视为固定的力心 $O'$，力心对单位质量的有心引力为 $\mu r$，$r$ 为力心到质量为 $m$ 的质点 $Q$ 的位矢大小，质点 $Q$ 被限制在圆周上运动。
\begin{enumerate}
  \item 求质点 $Q$ 从 $B$ 点由静止出发转过 $\phi$ 角有心力所做的功；
  \item 求质点通过第二象限所经历的时间。
\end{enumerate}

\item 判断力 $\vec{F} = x^{2}y^{2}\,\vec{i} + x^{2}y^{2}\,\vec{j}$ 是不是保守力？$\vec{F} = 2xy\,\vec{i} + x^{2}\,\vec{j}$ 呢？

\item 把一个物体从地球表面上沿铅垂方向以第二宇宙速度 $v_0 = \sqrt{\dfrac{2GM}{R_e}}$ 发射出去，阻力忽略不计。求物体从地面飞行到与地心相距 $nR_e$ 处所经历的时间。

\item 用弹簧连接两个木板 $m_1$、$m_2$，弹簧压缩 $x_0$。求给 $m_2$ 上加多大的压力能使 $m_1$ 离开桌面？

\item（书中例题 3.15，p.126）物体 $M$ 悬于弹簧上，弹簧的弹性系数为 $k$，弹簧的原长与圆环的半径相等。不计摩擦力。求物体自弹簧的原长无初速度地沿圆环滑至最低点 $B$ 时所获得的动能。

\end{enumerate}

\end{document}
